\chapter{Probability, hitting times and other criteria}\label{chap:probability}

This chapter studies P\'olya's function $P(\theta)=K_{i\theta/2}(2\pi)$, the building block of his approximation to $\Xi$, as a model in which every statement of the Thorin programme is a theorem with an explicit witness: Biane's Brownian hitting times, a Liouville operator whose Dirichlet spectrum gives the zeros, a GGC with the squared zeros as Thorin atoms whose clock is the hitting time of an explicit diffusion, and a Klein Gordon criterion. The de Bruijn Newman flow shows where the analogy with $\Xi$ stops. The exercises at the end also cover sections now placed elsewhere: the law of the clock of $\Xi$ (Section~\ref{ch:clocklaw}), its Brownian boundaries (Section~\ref{ch:boundaries}), Li's criterion (Section~\ref{ch:li}), the de Bruijn Newman flow (Section~\ref{ch:dbn}) and the Dirichlet beta function (Section~\ref{sec:dbeta:beta}).
